% Experimental plan: multi-agent verification of extraction outputs
% Ticket 0054 Phase A deliverable
%
% SUPERSEDED by tickets 0097 / 0099. Historical preregistration of the
% 5-mode post-hoc LLM verification protocol. No longer included in
% report.pdf (the \input line was removed from report.tex in ticket 0099).
% Ticket 0059 (PR #246) established that this protocol is a dead end
% (0–10% inter-verifier agreement). Kept as artifact only.

\section{Exp\'erience~5~: v\'erification multi-agents}\label{sec:exp5}

La section~\ref{sec:verification} d\'ecrit cinq modes de v\'erification o\`u
un seul agent (ou outil) annote chaque ligne du CSV avec des citations.
L'auto-v\'erification (\texttt{self}) souffre d'un biais de confirmation~:
le m\^eme mod\`ele valide ses propres fabrications. La v\'erification crois\'ee
(\texttt{cross}) utilise un mod\`ele diff\'erent mais reste un avis unique. Cette
exp\'erience teste si l'agr\'egation de \emph{plusieurs} agents v\'erificateurs
ind\'ependants am\'eliore la d\'etection d'hallucinations et la qualit\'e de
l'\'evidence, au-del\`a de ce qu'un v\'erificateur unique peut atteindre.

\subsection{Motivation}\label{sec:multi-verif-motivation}

Trois r\'esultats du pipeline existant motivent cette exp\'erience~:

\begin{enumerate}
\item \textbf{L'auto-consistance am\'eliore le rappel, pas la pr\'ecision.}
  L'analyse de self-consistency (section~\ref{sec:self-consistency}) montre que
  le vote union sur 3~runs augmente le rappel (plus de centrales d\'ecouvertes)
  mais n'\'elimine pas les hallucinations~: un nom fabriqu\'e pr\'esent dans 2~runs
  sur~3 passe le vote majoritaire. La consolidation multi-run est orthogonale \`a
  la v\'erification de provenance.

\item \textbf{La v\'erification crois\'ee d\'etecte plus d'hallucinations que
  l'auto-v\'erification.}
  L'hypoth\`ese~3 de la section~\ref{sec:verif-hypotheses} pr\'edit que le mode
  \texttt{cross} attribue davantage de scores~0 que le mode \texttt{self}.
  Si confirm\'ee, agr\'eger plusieurs v\'erificateurs crois\'es devrait amplifier
  cet avantage.

\item \textbf{Le module sourçage fournit un substrat de v\'erification.}
  Le module sourçage (ticket~0038) ajoute des colonnes \texttt{Source\_1} et
  \texttt{Source\_2} par ligne. Un agent v\'erificateur peut cibler chaque
  citation individuellement sans r\'e-extraire l'inventaire complet.
\end{enumerate}

\subsection{Choix du protocole}\label{sec:protocol-choice}

Quatre protocoles candidats ont \'et\'e \'evalu\'es contre l'infrastructure
existante~:

\begin{description}
\item[Vote majoritaire multi-run] D\'ej\`a impl\'ement\'e dans
  \texttt{self\_consistency.py}. Agr\`ege $k$~extractions du m\^eme mod\`ele par
  nom. Limite~: ne v\'erifie pas la \emph{provenance}, seulement la
  \emph{reproductibilit\'e}. Un nom hallucin\'e de mani\`ere reproductible passe le
  filtre.

\item[V\'erification crois\'ee multi-agents] Extension directe du mode
  \texttt{cross}~: $k$~mod\`eles v\'erificateurs ind\'ependants annotent chaque
  ligne, puis un vote majoritaire sur les scores d'\'evidence d\'etermine le
  verdict final. Compatible avec l'infrastructure existante~: chaque appel
  utilise \texttt{verify\_cross()} avec un mod\`ele diff\'erent.

\item[D\'ebat] Deux mod\`eles argumentent pour/contre chaque ligne disput\'ee,
  un juge tranche. Co\^ut \'elev\'e (3~appels LLM par ligne disput\'ee).
  N\'ecessite un prompt de d\'ebat sp\'ecifique non encore impl\'ement\'e.
  Report\'e \`a un travail futur.

\item[Ancrage dans les sources (\emph{source grounding})] Un agent
  v\'erificateur r\'ecup\`ere le document cit\'e dans \texttt{Source\_1} et v\'erifie
  que la donn\'ee revendiqu\'ee y figure. N\'ecessite l'acc\`es au corpus RAG
  complet et un m\'ecanisme de r\'ecup\'eration cibl\'ee. Faisable mais hors
  p\'erim\`etre de cette exp\'erience (requiert une int\'egration RAG-v\'erification
  non triviale).
\end{description}

\paragraph{Protocole retenu~: v\'erification crois\'ee multi-agents avec vote
majoritaire.}
Ce protocole maximise la r\'eutilisation de l'infrastructure existante
(\texttt{verify\_cross} appel\'e $k$~fois avec des mod\`eles diff\'erents) tout en
apportant la diversit\'e n\'ecessaire pour d\'epasser les angles morts d'un
v\'erificateur unique. Le vote majoritaire sur les scores d'\'evidence est le
m\'ecanisme d'agr\'egation le plus simple et le plus interpr\'etable.

\subsection{Sp\'ecification du protocole}\label{sec:multi-verif-spec}

\subsubsection{Panel de v\'erificateurs}\label{sec:verifier-panel}

Trois mod\`eles sont s\'electionn\'es pour maximiser la diversit\'e sur trois
axes~: fournisseur, architecture et base d'entra\^inement.

\begin{center}
\small
\begin{tabular}{llllrr}
\toprule
\textbf{Mod\`ele} & \textbf{Fournisseur} & \textbf{Pays} & \textbf{Arch.}
  & \textbf{In (\$/Mtok)} & \textbf{Out (\$/Mtok)} \\
\midrule
Claude Sonnet~4.6  & Anthropic & US & Dense   & 3,00 & 15,00 \\
GPT-5.4            & OpenAI    & US & Dense   & 2,50 & 10,00 \\
DeepSeek V3.2      & DeepSeek  & CN & MoE     & 0,26 &  0,38 \\
\bottomrule
\end{tabular}
\end{center}

\noindent Justification du panel~:
\begin{itemize}
\item \textbf{Diversit\'e de fournisseur}~: trois laboratoires distincts
  (Anthropic, OpenAI, DeepSeek) minimisent la corr\'elation des biais
  d'entra\^inement.
\item \textbf{Diversit\'e g\'eographique}~: deux mod\`eles US, un mod\`ele chinois.
  Pour l'extraction de donn\'ees vietnamiennes, les mod\`eles chinois ont
  montr\'e une meilleure connaissance des noms propres et des documents
  r\'egionaux (section~\ref{sec:census-results}).
\item \textbf{Diversit\'e architecturale}~: deux mod\`eles denses, un MoE.
  Les architectures MoE activent des experts diff\'erents selon le contexte,
  ce qui peut produire des erreurs d\'ecorr\'el\'ees.
\item \textbf{Co\^ut}~: DeepSeek V3.2 est ${\sim}40\times$ moins cher que
  Sonnet~4.6 en tokens de sortie. Le panel couvre un spectre de prix large
  pour mesurer le rendement marginal de chaque v\'erificateur suppl\'ementaire.
\end{itemize}

\subsubsection{Entr\'ees}\label{sec:multi-verif-inputs}

Les extractions de base proviennent des meilleures configurations identifi\'ees
dans les exp\'eriences pr\'ec\'edentes~:

\begin{enumerate}
\item \textbf{DeepSeek V3.2 d\'ecompos\'e} (F1 $\approx$ 98,8\,\%) --- meilleur
  score global, d\'ej\`a configur\'e dans \texttt{sweeps.verification}.
\item \textbf{Opus~4.6 multitour} (F1 $\approx$ 98,2\,\%, 169~centrales) ---
  mod\`ele frontier de r\'ef\'erence, haut rappel mais 6~hallucinations.
\item \textbf{GPT-5.4 RAG} --- repr\'esente la m\'ethode RAG wholesale.
\end{enumerate}

\noindent Pour chaque extraction de base, les colonnes \texttt{Source\_1} et
\texttt{Source\_2} sont pr\'esentes (issues du module sourçage ou ajout\'ees par
l'\'etape de v\'erification).

\subsubsection{Proc\'edure}\label{sec:multi-verif-procedure}

Pour chaque extraction de base~:

\begin{enumerate}
\item \textbf{V\'erification ind\'ependante.} Chaque v\'erificateur $V_i$
  ($i \in \{1, 2, 3\}$) re\c{c}oit le CSV extrait (noms, attributs, citations)
  et produit un score d'\'evidence par ligne (rubrique 0--4,
  section~\ref{sec:evidence-rubric}). Les v\'erificateurs ne voient ni le prompt
  d'extraction original ni les sorties des autres v\'erificateurs.

\item \textbf{Agr\'egation par vote majoritaire.} Pour chaque ligne, le score
  d'\'evidence final est la \emph{m\'ediane} des $k=3$ scores individuels. La
  m\'ediane est pr\'ef\'er\'ee \`a la moyenne car elle est robuste aux valeurs
  aberrantes (un v\'erificateur trop indulgent ou trop s\'ev\`ere).

\item \textbf{Classification finale.} Une ligne est class\'ee comme~:
  \begin{description}
  \item[V\'erifi\'ee] si le score m\'edian $\geq 3$ (au moins une source primaire
    selon la majorit\'e des v\'erificateurs).
  \item[Incertaine] si le score m\'edian $\in \{1, 2\}$ (aucune source primaire
    confirm\'ee par la majorit\'e).
  \item[Fabriqu\'ee] si le score m\'edian $= 0$ (hallucination d\'etect\'ee par la
    majorit\'e).
  \end{description}

\item \textbf{R\`egle de bris d'\'egalit\'e.} Avec $k=3$ v\'erificateurs et des
  scores entiers 0--4, la m\'ediane est toujours d\'efinie sans ambigu\"it\'e
  (\'el\'ement central du triplet tri\'e). Si $k$ est pair dans une extension
  future, prendre $\lfloor \text{m\'ediane} \rfloor$ (arrondir vers le bas =
  conservateur).
\end{enumerate}

\subsubsection{Conditions exp\'erimentales}\label{sec:multi-verif-conditions}

Le design factoriel croise 3~extractions de base $\times$ 4~conditions de
v\'erification~:

\begin{center}
\small
\begin{tabular}{clcr}
\toprule
\textbf{\#} & \textbf{Condition} & \textbf{$k$ v\'erificateurs} & \textbf{Appels LLM} \\
\midrule
C0 & Non v\'erifi\'e (baseline)           & 0 & 0 \\
C1 & Cross unique (Sonnet~4.6)          & 1 & 1 \\
C2 & Cross double (Sonnet + DeepSeek)   & 2 & 2 \\
C3 & Cross triple (Sonnet + DeepSeek + GPT-5.4) & 3 & 3 \\
\bottomrule
\end{tabular}
\end{center}

\noindent Les conditions C1 et C2 sont des sous-ensembles de C3, permettant de
mesurer le rendement marginal de chaque v\'erificateur suppl\'ementaire.
La condition C1 correspond au mode \texttt{cross} existant et sert de
contr\^ole.

\paragraph{R\'ep\'etitions.} Les modes LLM sont stochastiques. Chaque condition
$\times$ extraction est r\'ep\'et\'ee $n=3$~fois pour estimer la variance.
Le mode non v\'erifi\'e (C0) est d\'eterministe (1~run).

\paragraph{Total.} $3 \text{ bases} \times (1 + 3 \times 3) = 30$~conditions,
soit $3 + 27 = 30$~runs dont 27~appels LLM.

\subsection{Estimation des co\^uts}\label{sec:multi-verif-cost}

Chaque appel de v\'erification traite ${\sim}150$~lignes CSV. Le prompt de
v\'erification (liste de centrales + consigne de citation) consomme
${\sim}2\,000$~tokens en entr\'ee et ${\sim}4\,000$~tokens en sortie.

\begin{center}
\small
\begin{tabular}{lrrr}
\toprule
\textbf{V\'erificateur} & \textbf{Co\^ut/appel} & \textbf{Appels (3 bases $\times$ 3 reps)}
  & \textbf{Total} \\
\midrule
Sonnet~4.6   & \$0,066 & 9 & \$0,59 \\
GPT-5.4      & \$0,045 & 9 & \$0,41 \\
DeepSeek V3.2 & \$0,002 & 9 & \$0,02 \\
\midrule
\textbf{Total} & & 27 & \textbf{\$1,02} \\
\bottomrule
\end{tabular}
\end{center}

\noindent Co\^ut total estim\'e~: ${\sim}1$\,\$ pour 27~appels LLM. Le budget de
5\,\$ configur\'e dans \texttt{sweeps.verification} est largement suffisant.
DeepSeek V3.2 repr\'esente $<$2\,\% du co\^ut total malgr\'e un tiers des appels,
illustrant l'asym\'etrie de prix entre mod\`eles frontier.

\subsection{Analyse pr\'evue}\label{sec:multi-verif-analysis}

\paragraph{M\'etrique principale.} Pour chaque condition, la courbe
pr\'ecision-couverture (section~\ref{sec:prec-cov}) est calcul\'ee aux seuils
$t \in \{0, 1, 2, 3, 4\}$. La comparaison porte sur~:

\begin{itemize}
\item \textbf{Pr\'ecision au seuil $t=3$.} Fraction des centrales retenues qui
  sont dans le r\'ef\'erentiel, apr\`es filtrage des lignes dont le score m\'edian
  est $< 3$.
\item \textbf{Couverture au seuil $t=3$.} Fraction des centrales retenues par
  rapport au total avant filtrage.
\item \textbf{D\'etection d'hallucinations.} Nombre de centrales recevant un
  score m\'edian de~0 (fabriqu\'ees) dans chaque condition.
\item \textbf{Rendement marginal.} $\Delta$pr\'ecision et $\Delta$d\'etection
  entre C1, C2 et C3.
\end{itemize}

\paragraph{M\'etrique secondaire.} Le co\^ut cumul\'e par condition (en \$)
rapporte au gain de pr\'ecision pour \'evaluer le rapport co\^ut-efficacit\'e.

\subsection{Hypoth\`eses pr\'e-enregistr\'ees}\label{sec:multi-verif-hypotheses}

Les hypoth\`eses sont formul\'ees \emph{avant} l'ex\'ecution de l'exp\'erience.

\begin{enumerate}
\item \textbf{H1~: la v\'erification multi-agents d\'etecte plus d'hallucinations
  que la v\'erification crois\'ee unique.}
  Le nombre de lignes recevant un score m\'edian de~0 dans C3 (3~v\'erificateurs)
  est sup\'erieur ou \'egal \`a celui dans C1 (1~v\'erificateur). Justification~:
  chaque mod\`ele a des angles morts diff\'erents~; un mod\`ele peut confirmer par
  erreur une fabrication qu'un autre d\'etecte.

\item \textbf{H2~: le rendement marginal du troisi\`eme v\'erificateur est
  d\'ecroissant.}
  Le gain de pr\'ecision entre C2 et C3 est inf\'erieur au gain entre C1 et C2.
  Justification~: les deux premiers v\'erificateurs couvrent la majorit\'e des
  erreurs faciles~; le troisi\`eme n'apporte qu'un gain marginal sur les cas
  limites.

\item \textbf{H3~: la couverture au seuil $t=3$ diminue avec le nombre de
  v\'erificateurs.}
  Plus de v\'erificateurs signifie un filtre plus strict (la m\'ediane de
  3~scores est plus conservatrice que le score d'un seul v\'erificateur).
  La couverture de C3 est inf\'erieure \`a celle de C1 au seuil $t=3$.

\item \textbf{H4~: DeepSeek V3.2 est le v\'erificateur le plus co\^ut-efficace.}
  Le rapport $\Delta$hallucinations d\'etect\'ees / co\^ut est le plus \'elev\'e pour
  DeepSeek V3.2. Justification~: \`a ${\sim}40\times$ moins cher que Sonnet~4.6,
  m\^eme une performance de v\'erification mod\'er\'ee produit un meilleur ratio.
\end{enumerate}

\paragraph{Seuils de falsification.} Une hypoth\`ese directionnelle est
r\'efut\'ee si l'intervalle de confiance bootstrap \`a 95\,\% sur la diff\'erence
exclut z\'ero dans la direction oppos\'ee \`a la pr\'ediction. H4 est \'evalu\'ee
par le ratio co\^ut-efficacit\'e (hallucinations d\'etect\'ees par dollar), sans
seuil absolu.

\subsection{Impl\'ementation}\label{sec:multi-verif-implementation}

L'impl\'ementation requiert deux modifications de \texttt{query\_verification.py}~:

\begin{enumerate}
\item \textbf{Mode \texttt{multi\_cross}.} Nouveau mode qui appelle
  \texttt{verify\_cross()} s\'equentiellement pour chaque mod\`ele du panel, puis
  agr\`ege les scores par m\'ediane. La sortie est un CSV annot\'e avec les
  colonnes~: \texttt{evidence\_score\_v1}, \texttt{evidence\_score\_v2},
  \texttt{evidence\_score\_v3}, \texttt{evidence\_score\_median}.

\item \textbf{Configuration dans \texttt{experiments.toml}.} Un nouveau sweep
  \texttt{verification\_multi} sp\'ecifie le panel de v\'erificateurs et les
  extractions de base.
\end{enumerate}

\noindent Ces modifications constituent le livrable du ticket~0055 (Phase~B,
bloqu\'e par ce ticket).

\subsection{Limites du design}\label{sec:multi-verif-limits}

\begin{itemize}
\item \textbf{Taille du panel.} $k=3$ v\'erificateurs est un minimum pour le
  vote majoritaire. Un panel de 5 ou 7 augmenterait la robustesse mais
  multiplie le co\^ut proportionnellement.

\item \textbf{M\^eme prompt de v\'erification.} Tous les v\'erificateurs re\c{c}oivent
  le m\^eme prompt (\texttt{\_VERIFICATION\_PROMPT}). La diversit\'e provient des
  mod\`eles, pas des prompts. Un design futur pourrait varier aussi les
  instructions de v\'erification.

\item \textbf{Pas d'acc\`es aux documents sources.} Les v\'erificateurs citent des
  sources de m\'emoire, sans acc\`es au corpus RAG. Le mode \texttt{web} reste
  sup\'erieur pour la v\'erification factuelle. Ce protocole mesure la
  \emph{convergence inter-mod\`eles}, pas la v\'erit\'e terrain documentaire.

\item \textbf{T\^ache unique.} Le design est test\'e sur les centrales thermiques
  du Vietnam uniquement. La g\'en\'eralisabilit\'e \`a d'autres pays ou classes
  de composants n'est pas \'etablie.
\end{itemize}

\subsection{Lien avec les autres exp\'eriences}\label{sec:multi-verif-links}

Cette exp\'erience s'inscrit dans la progression m\'ethodologique du pipeline~:

\begin{enumerate}
\item \textbf{Exp\'erience~1 (census)}~: extraction brute, pas de v\'erification.
\item \textbf{Exp\'erience~2 (RAG)}~: extraction avec contexte documentaire.
\item \textbf{Exp\'erience~3 (ablation)}~: optimisation du prompt, module sourçage.
\item \textbf{Exp\'erience~4 (v\'erification 5~modes)}~: un v\'erificateur par
  mode, courbe pr\'ecision-couverture.
\item \textbf{Exp\'erience~5 (v\'erification multi-agents)}~: panel de
  v\'erificateurs, rendement marginal de la diversit\'e.
\end{enumerate}

\noindent L'exp\'erience~5 r\'epond \`a la question~: \emph{au-del\`a d'un seul
v\'erificateur crois\'e, la diversit\'e inter-mod\`eles am\'eliore-t-elle la
d\'etection d'hallucinations, et \`a quel co\^ut~?}
